% Staged replacement for the opening prose of article_content.tex.
Writing reaches a reader as visible marks. Script, typeface, handwriting and
damage change those marks without necessarily changing their interpretation.
Conversely, historical or unfamiliar forms may have no agreed character
identity in a chosen encoding or tokenizer. A visual language interface would
retain the marks while learning to interpret and continue them. Its intended
output is readable writing, including explanatory passages and, eventually,
attested forms absent from the model's preparation fonts. This motivation does
not establish that image representations carry more usable information per
unit of computation or reproduce human reading mechanisms.

Pixel-space language generation has established precedents. PIXAR studies
autoregressive text generation from pixel representations and identifies
readability as a central difficulty\cite{pixar}. PixelRNN and conditional
PixelCNN provide normalized image-density formulations\cite{pixelrnn,pixelcnn}.
Their existence motivates a trainable visual writing distribution, but neither
an image likelihood nor a sharp glyph establishes that an answer follows its
question. Replacing a character inventory with physical image variables also
does not remove discrete computation: a binary raster and a categorical
distribution over small pixel masks are discrete representations of appearance.

We study the coupling directly. A prompt image supplies context; a neural
writer emits a glyph image; the model rereads the actual emitted pixels before
deciding what to write next and when to stop. Typed questions may be rendered
at the input boundary, and transcription may follow completed image generation.
Neither operation supplies answer identities to the deployed writing loop.
Offline text, teachers and pretrained foundations remain permitted sources of
supervision and inherited knowledge. Independent inference therefore differs
from learning a language entirely from images without textual pretraining.

The initial language setting uses Chinese writing in one canonical font and
ordered glyph cells. This isolates conditional shape generation before raw
page reading, handwriting or historical-script transfer. Earlier compact
complete-sequence fits produce zero strict complete answers at their fixed
checkpoints. The subsequent full continuation trial recomputes all 36 answer
layers while preserving frozen prompt memory, exposing complete QA and book
targets rather than only isolated first glyphs. Its five prescribed checkpoints,
paired questions, erased inputs, long continuations and retained successes
separate physical fidelity from complete language. The final model improves
some short answers relative to its initial checkpoint but fails the pilot.

The study contributes an explicit emitted-pixel inference interface, controlled
complete-language comparisons and evidence that supplied-history writing and
one-GPU execution are insufficient qualifications. The failed full trial
motivates a separate conditional joint glyph writer. That head-only fit is
complete and fails its physical prerequisite; its original review credits
are preserved with detected blank-flag errors rather than treated as corrected
physical accuracy. The probability diagnosis in turn motivates training the
continuation state through complete-image likelihood. The exact state-gradient
implementation passes a real-source computational preflight, but its completed
fit and language remain unmeasured. The larger objective is useful
visual language across written forms, with historical and unencoded writing
assessed on its own provenance, shape and meaning. That objective remains open.
